\chapter{Không phải đứa nào cũng tự tin sẵn}
\markboth{Không phải đứa nào cũng tự tin sẵn}{Không phải đứa nào cũng tự tin sẵn}

\section{Em không ngu, em chỉ vô bài chậm}
\begin{truyen}{Em không ngu, em chỉ vô bài chậm}{Classroom Talk}
\chuhoa{H}{ọc sinh: ``Thầy đừng nhìn em kiểu bất lực. Em không ngu. Em chỉ vô bài chậm hơn tụi kia thôi.''}
\textit{Student: ``Please do not look at me like I am hopeless. I am not stupid. I just pick things up more slowly than the others.''}\\[4pt]

\teacherpair{Nói vậy là đúng hơn nhiều đứa chỉ biết tự chửi mình. Vào chậm không đáng sợ bằng tưởng mình hết cứu.}{That is already more accurate than the way many students only insult themselves. Learning slowly is not as dangerous as believing you are beyond help.}

\studentpair{Nhưng chậm riết em cũng nản.}{But when it stays slow for too long, I get discouraged.}

\teacherpair{Vậy mình tìm cách học hợp với nhịp của em, chứ không lấy nhịp người khác đập lên đầu em nữa.}{Then we need to find a way of learning that fits your pace instead of using somebody else’s pace to hit you over the head.}
\end{truyen}

\begin{baihoc}
Tự nhận mình chậm còn sửa được. Tự kết án mình ngu là đang tự khóa cửa.
\textit{(Admitting you are slower can be worked with. Labeling yourself stupid shuts the door.)}
\end{baihoc}

\begin{ghinhoanh}
\textbf{Short English to remember:}

Slow is not hopeless.

Wrong labels damage real progress.
\end{ghinhoanh}

\begin{tuhocnhanh}
1. Em đang thiếu khả năng hay thiếu cách học hợp mình? \textit{(Are you lacking ability, or the right learning method?)}

2. Nhịp học nào khiến em đỡ bị ngợp nhất? \textit{(What learning pace helps you feel less overwhelmed?)}
\end{tuhocnhanh}
\ngancach

\section{Bạn kia có nền sẵn rồi, sao so với em được}
\begin{truyen}{Bạn kia có nền sẵn rồi, sao so với em được}{Classroom Talk}
\chuhoa{H}{ọc sinh: ``Bạn đó học thêm từ năm nào rồi, về nhà còn có người kèm. Còn em tự bơi. Sao cứ lôi em ra so?''}
\textit{Student: ``That person has had tutoring for years and someone still helps at home. I am on my own. So why do people keep comparing me with them?''}\\[4pt]

\teacherpair{Em nói đúng một nửa. Xuất phát điểm khác nhau thật. Nhưng nếu cứ giữ nửa đúng đó để khỏi phải cố, em sẽ tự thua luôn nửa còn lại.}{You are right about half of it. Starting points are truly different. But if you hold on to that half-truth just to avoid trying, you will lose the other half as well.}

\studentpair{Nửa còn lại là gì?}{And what is the other half?}

\teacherpair{Là dù chậm hơn, em vẫn phải chạy phần đường của em. Không ai chạy hộ.}{The other half is this: even if you are slower, you still have to run your own part of the road. No one can run it for you.}
\end{truyen}

\begin{baihoc}
Đừng phủ nhận bất công. Nhưng cũng đừng dùng bất công làm cớ để đứng yên.
\textit{(Do not deny unfairness, but do not let it become a reason to remain still.)}
\end{baihoc}

\begin{ghinhoanh}
\textbf{Short English to remember:}

Different starts are real.

Your path is still yours to run.
\end{ghinhoanh}

\begin{tuhocnhanh}
1. Em hay than về xuất phát điểm đến mức nào? \textit{(How often do you dwell on your starting point?)}

2. Phần việc nào vẫn là trách nhiệm riêng của em? \textit{(What part of the work is still your responsibility?)}
\end{tuhocnhanh}
\ngancach

\section{Em cứ bị đơ khi lên bảng}
\begin{truyen}{Em cứ bị đơ khi lên bảng}{Classroom Talk}
\chuhoa{H}{ọc sinh: ``Ngồi dưới em còn nghĩ được. Lên bảng là đầu em trắng trơn. Nhìn cả lớp là em đơ.''}
\textit{Student: ``When I sit below, I can still think. But once I go to the board, my mind goes blank. The moment everyone looks at me, I freeze.''}\\[4pt]

\teacherpair{Em sợ sai hay sợ bị nhìn thấy mình sai?}{Are you afraid of being wrong, or afraid that people will see you being wrong?}

\studentpair{Chắc sợ cái thứ hai nhiều hơn.}{Probably the second one more than the first.}

\teacherpair{Vậy vấn đề không chỉ là kiến thức. Em còn đang bị sĩ diện bóp cổ. Mình phải tập chịu quê một chút mới khá lên được.}{Then the problem is not only knowledge. Your pride is choking you as well. You need to practice tolerating a little embarrassment if you want to improve.}
\end{truyen}

\begin{baihoc}
Muốn đỡ run, không chỉ học bài. Còn phải học chịu cảm giác mình chưa giỏi trước mặt người khác.
\textit{(To become less anxious, you must train not only knowledge but also your tolerance for visible imperfection.)}
\end{baihoc}

\begin{ghinhoanh}
\textbf{Short English to remember:}

Fear of looking foolish is powerful.

Growth needs some public discomfort.
\end{ghinhoanh}

\begin{tuhocnhanh}
1. Em sợ sai hay sợ mất mặt hơn? \textit{(Do you fear being wrong, or losing face more?)}

2. Em có thể tập nói trước một nhóm nhỏ trước không? \textit{(Can you practice speaking before a smaller group first?)}
\end{tuhocnhanh}
\ngancach

\section{Cứ bị nói lười riết em cũng chán cố}
\begin{truyen}{Cứ bị nói lười riết em cũng chán cố}{Classroom Talk}
\chuhoa{H}{ọc sinh: ``Em biết em chưa tốt. Nhưng nghe bị gắn mác lười suốt ngày, riết em cũng muốn mặc kệ luôn.''}
\textit{Student: ``I know I am not doing well. But hearing myself labeled lazy all the time makes me want to give up altogether.''}\\[4pt]

\teacherpair{Người lớn nhiều khi nóng miệng. Nhưng em cũng phải tự hỏi: mình đang bị oan hoàn toàn, hay mình có phần để người ta nói?}{Adults often speak too harshly. But you also need to ask yourself this: are you completely being treated unfairly, or is there some truth in what they say?}

\studentpair{Có phần thật. Nhưng nghe hoài khó chịu lắm.}{There is some truth in it. But hearing it over and over still feels awful.}

\teacherpair{Khó chịu là thật. Nên mình cần hai việc: bớt cái khiến người ta nói, và đừng lấy lời nói đó làm lý do hỏng tiếp.}{That discomfort is real. So we need two things: reduce the behavior that makes people say it, and stop using their words as an excuse to keep failing.}
\end{truyen}

\begin{baihoc}
Không phải lời chê nào cũng đúng. Nhưng nếu lời chê làm mình buông luôn, phần thiệt lại về mình.
\textit{(Not every criticism is fair, but if it makes you quit, you still lose.)}
\end{baihoc}

\begin{ghinhoanh}
\textbf{Short English to remember:}

Labels hurt.

Do not let them become your excuse.
\end{ghinhoanh}

\begin{tuhocnhanh}
1. Em bị dán nhãn gì nhiều nhất? \textit{(What label do people place on you most often?)}

2. Em cần sửa điều gì thật sự để tự gỡ nhãn đó? \textit{(What real change would help remove that label?)}
\end{tuhocnhanh}
\ngancach